% ====================================================================
% FILE: content/16_modules_master.tex
% Архитектура модулей и перекрестные структуры по уровням
% ====================================================================

\section{Архитектура модулей и перекрестные структуры по уровням}
\label{sec:modules_master}

Этот раздел даёт обзор высокого уровня расширенных структурных модулей
Онтологии континуумов. Каждый модуль разработан в техническом атласе;
здесь мы кратко суммируем его роль и связи, чтобы основной стержень имел
целостную связь между формальной доктриной и разделами закрытия Core~1.3.3.

Расширенные модули:

\begin{itemize}
\item \textbf{Мастер уровней K} (раздел~the K-level overview chapter):
  глобальный обзор всех континуумов \(K_0\)–\(K_{12}\), их пространств
  состояний, осей, порогов и ключевых процессов.

\item \textbf{Пространства вложения / M-пространства}
  (раздел~the M-space overview chapter):
  определение метапространств \(M_x\), которые ограничивают и параметризуют
  эволюцию континуумов \(K_x\).

\item \textbf{Cross-level / cross-K структуры} (раздел~the cross-level structure chapter):
  операторы, отображения и ландшафтные структуры, связывающие различные
  уровни иерархии и генерирующие фазовые переходы между ними.

\item \textbf{Процессы} (раздел~the process taxonomy chapter):
  унифицированное описание операторов эволюции
  \(E = (F,G,H,Q,R,S,U)\) на разных уровнях, включая динамику напряжения,
  пороги, потоки и циклы.

\item \textbf{Циклы} (раздел~the cycle taxonomy chapter):
  структурные и динамические циклы \(C_j\) как носители непрерывности,
  памяти и временной организации.

\item \textbf{Предсказания} (раздел~the bounded prediction-route chapter):
  предсказания на уровне и между уровнями, эмпирические сигнатуры и
  качественные режимы, вытекающие из формализма.

\item \textbf{Эксперименты} (раздел~the experiment-design chapter):
  концептуальные и прикладные экспериментальные дизайны для исследования
  континуумов, переходов и пороговых ландшафтов.

\item \textbf{Структура фальсифицируемости} (раздел~the falsifiability chapter):
  явные условия, при которых Онтология континуумов может быть опровергнута,
  ограничена или вынуждена перестроить ядро.

\item \textbf{Jets / структура локальной эволюции} (раздел~the local-evolution chapter):
  локальные разложения оператора эволюции и краткосрочной динамики
  около заданных конфигураций \(K_x\) и \(M_x\).
\end{itemize}

Совместно эти модули уточняют и делают рабочей основу ядра на всех уровнях
приложений, сохраняя при этом единую согласованную иерархию континуумов и
метапространств.
